\subsection{Algoritmo}

El algoritmo crea una matriz de adyacencia para representar un grafo, donde cada entrada $cost[i][j]$ almacena el costo de viajar desde el vértice $i$ al vértice $j$. 

Inicialmente, todos los costos son cero. El usuario ingresa los costos de las aristas, y el programa los almacena en la matriz. Si un costo no se ingresa, se asume un valor predeterminado de 15000 (considerado infinito en este contexto). 

Luego, la matriz se ajusta para que sea simétrica (el costo de $i$ a $j$ es el mismo que de $j$ a $i$), lo que indica que el grafo es no dirigido.

El algoritmo luego solicita un vértice de salida $v$. A partir de este vértice, calcula las distancias más cortas a todos los demás vértices utilizando una variante del algoritmo de Dijkstra. 

Mantiene un arreglo $dist$ donde $dist[i]$ almacena la distancia más corta conocida desde $v$ hasta $i$, y un arreglo $sol$ (solución) que marca qué vértices ya han sido visitados o "fijados" en el camino más corto. 

Inicialmente, $dist[v]$ se establece en 0 y todos los demás en infinito.

Enunciamos que el algoritmo se parece al algoritmo de Dijkstra porque\dots

Las estructuras de datos que utiliza son una matriz de adyacencia y dos arreglos. La matriz de adyacencia se utiliza para almacenar los costos de las aristas, y los arreglos se utilizan para almacenar las distancias más cortas y los vértices visitados.

Lo consideramos un algoritmo greedy porque\dots

La complejidad temporal del algoritmo es $O(n^2)$ y la complejidad espacial es $O(n^2)$. Porque \dots
